\begin{figure}[!htbp]
\centering
\resizebox{0.96\textwidth}{!}{%
\begin{tikzpicture}[
  card/.style={draw=accent!70, fill=blue!5, rounded corners=6pt,
    text width=3.6cm, minimum height=1.35cm, align=center,
    font=\footnotesize, inner sep=7pt},
  artifact/.style={draw=green!55!black, fill=green!8, rounded corners=6pt,
    text width=4.2cm, minimum height=1.35cm, align=center,
    font=\footnotesize\bfseries, inner sep=7pt},
  arr/.style={->, >=Stealth, thick, accent!75, shorten <=3pt, shorten >=3pt}
]
\node[card] (problem) at (0,0) {\textbf{Real problem}\\Dynamic evidence is missed by static crawling};
\node[card] (objectives) at (4.25,0) {\textbf{Clear objectives}\\Classify, interact, capture, trace, review};
\node[artifact] (artifact) at (8.75,0) {Designed artifact\\Open Web Catcher};
\node[card] (evaluation) at (13.25,0) {\textbf{Evaluation}\\Evidence quality, traceability, cost, failures};

\draw[arr] (problem) -- (objectives);
\draw[arr] (objectives) -- (artifact);
\draw[arr] (artifact) -- (evaluation);

\node[draw=gray!55, fill=gray!7, rounded corners=5pt, text width=12.6cm,
      minimum height=0.78cm, align=center, font=\footnotesize] at (6.65,-1.65)
{\textbf{Methodological decision:} DSRM is used as the sole backbone because the PFE contribution is a designed, demonstrated, and evaluated information-system artifact.};
\end{tikzpicture}
}
\caption{DSRM selection logic for an artifact-centered PFE project.}
\label{fig:methodology-selection}
\end{figure}
